% Compile with XeLaTeX twice. Requires Noto Sans CJK SC or FandolSong.
\documentclass[11pt,a4paper]{article}
\usepackage[margin=21mm,top=20mm,bottom=21mm,headheight=16pt]{geometry}
\usepackage{fontspec}
\IfFontExistsTF{Noto Sans CJK SC}
  {\setmainfont{Noto Sans CJK SC}[AutoFakeBold=2,AutoFakeSlant=0.15]}
  {\setmainfont{FandolSong-Regular}[AutoFakeBold=2,AutoFakeSlant=0.15]}
\XeTeXlinebreaklocale "zh"
\XeTeXlinebreakskip=0pt plus 1pt
\usepackage{amsmath,amssymb,mathtools}
\usepackage{xcolor,enumitem,booktabs,tabularx,fancyhdr,lastpage}
\usepackage{hyperref}
\definecolor{ink}{HTML}{18394B}
\definecolor{pale}{HTML}{EDF3F5}
\hypersetup{colorlinks=true,urlcolor=ink,linkcolor=ink,
 pdftitle={C01 生成泛函、源与关联函数},pdfauthor={},
 pdfsubject={欧氏路径积分中的均值、涨落、响应与全息字典}}
\pagestyle{fancy}\fancyhf{}
\fancyhead[L]{\small\color{ink}C01\quad 生成泛函、源与关联函数}
\fancyhead[R]{\small 2026-11-03}
\fancyfoot[L]{\small 理论物理阅读材料}
\fancyfoot[R]{\small\thepage\ /\ \pageref*{LastPage}}
\setlength{\parindent}{0pt}
\setlength{\parskip}{0.45em}
\linespread{1.12}
\setlist{nosep,leftmargin=1.7em,topsep=0.35em}
\setlength{\fboxsep}{7pt}
\newcommand{\dd}{\mathrm d}
\newcommand{\OO}{\mathcal O}
\newcommand{\ZZ}{\mathcal Z}
\newcommand{\intx}{\int\dd^d x\,}
\newcommand{\avg}[1]{\left\langle #1\right\rangle}
\newcommand{\physicsnote}[1]{\par\smallskip\noindent\colorbox{pale}{\parbox{\dimexpr\linewidth-2\fboxsep\relax}{\textbf{物理解释}\quad #1}}\par\smallskip}
\newcommand{\question}[1]{\textbf{本节要回答：}#1\par}
\newcommand{\step}[1]{\textbf{#1}\par}
\begin{document}
\begin{center}
{\LARGE\bfseries\color{ink}C01：生成泛函、源与关联函数}\\[0.55em]
{\large 从改变构型权重，到读取均值、涨落与响应}\\[0.35em]
{\small 2026-11-03\quad 主线含练习约 75--100 分钟；第 8 节选读约 15--20 分钟}
\end{center}
\textbf{本次目标。}假定已经见过路径积分，不从时间切片重新开始。重点是亲手推导
\(Z[J]\)、\(W[J]\) 的源导数，理解“连通”为什么出现，并把这套语言接到全息中的边界源与响应。第 9 节用于检验理解，不要求一次完成所有延伸内容。

\section{源究竟改变了什么？}
\question{为什么先人为加入一个 \(J\)，再对它求导？}
采用欧氏形式、\(\hbar=1\)，并先考虑一个实的玻色标量算符 \(\OO[\phi](x)\)。约定
\begin{equation}\label{eq:definitions}
 \ZZ[J]=\int\mathcal D\phi\,
 e^{-S_E[\phi]+\intx J(x)\OO(x)},\qquad
 Z[J]=\frac{\ZZ[J]}{\ZZ[0]},\qquad W[J]=\ln Z[J].
\end{equation}
这里 \(x\) 包括欧氏时间；\(\intx\) 是整个欧氏时空的积分。因而 \(Z[0]=1\)、\(W[0]=0\)。以下把状态、边界条件与正规化方式固定；\(J\) 只通过写出的线性耦合进入，\(\OO\) 和测度不另含 \(J\)。

\(J(x)\) 是\textbf{外部指定的函数}，不是本次路径积分中的积分变量。每给定一个 \(J\)，就得到一个不同的构型权重。例如，在实的正权重情形，正的局部 \(J\) 会相对提高该处 \(\OO\) 较大构型的权重。改变 \(J\) 后重新平均，便能探查系统愿意怎样改变。

若 \(\OO\) 是磁化密度，\(J\) 可以是磁场；哈密顿量扰动写为
\(H_J=H_0-\int\dd^{d-1}\mathbf x\,J\OO\)，对应欧氏指数中的 \(+\int J\OO\)。对于其他算符，\(J\) 也可以只是计算用的探针，不一定能在实验中独立调节。

\physicsnote{生成泛函记录的是：对各种源分布施加扰动时，路径积分如何变化。它把逐个计算的关联函数，组织成同一个泛函对外部参数的各阶导数。\(Z[0]=1\) 只是归一化选择，不表示无源系统没有能量或涨落。}
\textbf{符号提醒。}本文始终使用欧氏 \(W=\ln Z\) 与指数中的 \(+J\OO\)。Lorentzian 文献常用 \(W=-i\ln Z\)，不能直接搬用其中的 \(i\) 和符号。在质量量纲下，若 \([\OO]=\Delta\)，则 \([J]=d-\Delta\)，使 \(\int J\OO\) 无量纲。

\newpage
\section{一次源微分，为什么插入一个算符？}
\question{“对源求导”怎样从路径积分中取出某一点的 \(\OO\)？}
先把泛函微分当作多元微分的连续版本。定义
\begin{equation}
 \delta F[J]=\int\dd^d x\,\frac{\delta F[J]}{\delta J(x)}\,\delta J(x),
 \qquad \frac{\delta J(y)}{\delta J(x)}=\delta^{(d)}(y-x).
\end{equation}
因此 \(\delta(\int J\OO)/\delta J(x)=\OO(x)\)。这一步与普通公式
\(\partial e^{\sum_i J_i O_i}/\partial J_k=O_k e^{\sum_iJ_iO_i}\) 完全对应。

在固定源下，\textbf{归一化的}期望值定义为
\begin{equation}\label{eq:expectation}
 \avg{X}_J=\frac{1}{\ZZ[J]}\int\mathcal D\phi\,X[\phi]\,
 e^{-S_E[\phi]+\int J\OO}.
\end{equation}
对式 \eqref{eq:definitions} 求导，分母 \(\ZZ[0]\) 是常数，于是
\begin{align}
 \frac{\delta Z[J]}{\delta J(x)}
 &=\frac{1}{\ZZ[0]}\int\mathcal D\phi\,\OO(x)e^{-S_E+\int J\OO}\notag\\
 &=\frac{\ZZ[J]}{\ZZ[0]}\avg{\OO(x)}_J
 = Z[J]\avg{\OO(x)}_J.\label{eq:firstZ}
\end{align}
\textbf{“插入”}就是在对所有场构型求和时，多乘上一个 \(\OO(x)\)。它不是在积分结束后再向某条经典轨迹加一个粒子。插入哪个算符，由源耦合的是哪个算符决定。

重复求导得到
\begin{equation}\label{eq:nZ}
 \frac{1}{Z[J]}\frac{\delta^n Z[J]}{\delta J(x_1)\cdots\delta J(x_n)}
 =\avg{\OO(x_1)\cdots\OO(x_n)}_J.
\end{equation}
在 \(J=0\) 时，由 \(Z[0]=1\) 可以省去左侧的 \(1/Z\)。在非零源下则不能随手省略。这些是欧氏关联函数；在量子理论的算符表述中，对应适当的欧氏时间排序。

\physicsnote{先保留任意 \(J\)，是为了拥有一组可独立微分的探针；求导后再置 \(J=0\)，读到的就是原理论的关联。若一开始把 \(J\) 丢掉，只剩一个数 \(Z[0]\)，就失去了这些导数信息。}

\step{把“生成”看得更具体}
形式上，在零源附近展开
\begin{equation}
 Z[J]=1+\intx J(x)\avg{\OO(x)}_0
 +\frac12\int\dd^d x\dd^d y\,J(x)J(y)\avg{\OO(x)\OO(y)}_0+\cdots.
\end{equation}
关联函数是这个泛函展开的系数。连续场论中这些系数通常是分布；重合点需要正规化与重整化，不能把所有位置都当作可直接代入的普通函数值。

\newpage
\section{为什么取对数？连通项从哪里来？}
\question{为什么系统的平均值变化，与“减去平均值乘积”的两点函数有关？}
由式 \eqref{eq:firstZ}，
\begin{equation}
 \frac{\delta W}{\delta J(x)}=\frac{1}{Z}\frac{\delta Z}{\delta J(x)}
 =\avg{\OO(x)}_J.
\end{equation}
再求一次导数，关键是\textbf{归一化因子本身也依赖源}。设
\(N_x[J]=\int\mathcal D\phi\,\OO(x)e^{-S_E+\int J\OO}\)，则
\begin{align}
 \frac{\delta\avg{\OO(x)}_J}{\delta J(y)}
 &=\frac{1}{\ZZ}\frac{\delta N_x}{\delta J(y)}
   -\frac{N_x}{\ZZ^2}\frac{\delta\ZZ}{\delta J(y)}\notag\\
 &=\avg{\OO(x)\OO(y)}_J
   -\avg{\OO(x)}_J\avg{\OO(y)}_J\notag\\
 &\equiv G_{c,J}(x,y)
 =\frac{\delta^2 W}{\delta J(x)\delta J(y)}.\label{eq:covariance}
\end{align}
第二项直接来自归一化分母的求导；在正权重情形，可以理解为“改变权重后，总概率仍须归一化”。记
\(\Delta\OO(x)=\OO(x)-\avg{\OO(x)}_J\)，便有
\(G_{c,J}(x,y)=\avg{\Delta\OO(x)\Delta\OO(y)}_J\)。

\physicsnote{若两个位置的涨落同向出现，协方差通常为正；反向出现则可为负。单纯的 \(\avg{\OO(x)}\avg{\OO(y)}\) 只记录各自的背景均值，不描述共同涨落。因此二阶导数保留的是连通两点函数。零协方差一般不等于统计独立，不能把“没有两点关联”当作“没有任何关联”。}

\step{更高阶：对数把独立块拆开}
\(W\) 的各阶源导数是连通关联函数，也称累积量（cumulants）。直观上，若两个真正独立的子系统有
\(Z_{A+B}=Z_AZ_B\)，则 \(W_{A+B}=W_A+W_B\)；跨越两个子系统的混合导数自动为零。一般的组合关系则由 \(Z=e^W\) 重建所有连通块的乘积。

例如，\textbf{零源、零均值的高斯场}有
\begin{equation}
 \avg{\phi_1\phi_2\phi_3\phi_4}_0
 =G_{12}G_{34}+G_{13}G_{24}+G_{14}G_{23},\qquad
 \avg{\phi_1\phi_2\phi_3\phi_4}_{c,0}=0.
\end{equation}
其中 \(G_{ij}=\avg{\phi(x_i)\phi(x_j)}_0\)。完整四点函数可以非零，只是没有独立于两点函数的新四点连通信息。非零均值时，完整四点函数还含均值贡献，不能照抄这个三项式。

\newpage
\section{算到底：一个高斯变量的均值、涨落与响应}
\question{上述公式是否真的描述“外场推移平均值”的过程？}
先暂时不考虑时空，只计算普通积分。令 \(a>0\)，
\begin{equation}
 \ZZ(J)=\int_{-\infty}^{\infty}\dd q\,e^{-a q^2/2+Jq}.
\end{equation}
\(a\) 控制偏离原点的代价；源 \(J\) 给分布一个倾向。完成平方：
\begin{equation}
 -\frac a2q^2+Jq=-\frac a2\left(q-\frac Ja\right)^2+\frac{J^2}{2a}.
\end{equation}
将积分变量平移为 \(\eta=q-J/a\)，积分的宽度没有改变，得到
\begin{equation}
 \ZZ(J)=\sqrt{\frac{2\pi}{a}}\,e^{J^2/(2a)},\qquad
 Z(J)=e^{J^2/(2a)},\qquad W(J)=\frac{J^2}{2a}.
\end{equation}
于是均值、完整二阶矩和方差分别为
\begin{align}
 \bar q=\avg q_J&=W'(J)=\frac Ja,\notag\\
 \avg{q^2}_J&=\frac{Z''(J)}{Z(J)}=\frac1a+\frac{J^2}{a^2},\notag\\
 \avg{(q-\bar q)^2}_J&=W''(J)=\frac1a
 =\frac{\dd\bar q}{\dd J}.\label{eq:gaussianResponse}
\end{align}
这里的微分是\textbf{比较两个不同源下的分布}，没有涉及系统随时间怎样运动。只有均值被平移，围绕均值的涨落宽度保持不变，这是线性源耦合的高斯模型的特殊性质；相互作用理论的方差一般会随 \(J\) 改变。

\physicsnote{更大的 \(a\) 让分布更窄，也让同样的源更难推动均值，所以涨落与响应同时减小。\(J=0\) 时均值为零，方差却为 \(1/a\)：平均值为零绝不是每个构型都等于零。高斯理论也有真实的涨落。}

\step{为什么刚才的等式像“涨落决定响应”？}
在这个静态模型里，较大的自然涨落意味着权重更容易被小源偏置。若把指数改写为
\(-\beta kq^2/2+\beta f q\)，则 \(a=\beta k\)、\(J=\beta f\)。对\textbf{物理外力}求导得到
\begin{equation}
 \frac{\dd\bar q}{\dd f}=\beta\,\mathrm{Var}(q)=\frac1k.
\end{equation}
这说明温度因子是否出现，取决于源的定义；不能把对 \(J\) 与对 \(f\) 的导数混为一谈。此例也解释了为什么使用欧氏约定时，要先明确哪些因子已吸收到源里。

\newpage
\section{推广到场：为什么两点函数是算符的逆？}
\question{高斯积分的 \(1/a\)，在场论中变成了什么？}
考虑线性源耦合到实标量场 \(\OO=\phi\)，且
\begin{equation}
 S_E[\phi]=\frac12\intx\phi(x)A\phi(x),\qquad A=-\partial_E^2+m^2.
\end{equation}
这一写法已做分部积分；这里取合适的齐次边界条件或周期条件，使边界项消失。先在有限体积、有 UV 正规化的情形推导，要求 \(A\) 对称、正定、可逆。无质量理论的零模需另外处理，不能直接写成 \(A^{-1}\)。

定义 Green 核 \(G\) 为
\begin{equation}
 A_xG(x,y)=\delta^{(d)}(x-y),\qquad
 (GJ)(x)=\int\dd^d y\,G(x,y)J(y).
\end{equation}
\(A_x\) 作用在第一个位置上。指定微分表达式还不够，\textbf{求逆时必须指定边界条件}。不同的边界或状态问题，可能选择不同的 Green 函数。

用 \(\phi=\eta+GJ\) 完成泛函版本的平方：
\begin{equation}
 -\tfrac12\phi A\phi+J\phi
 =-\tfrac12\eta A\eta+\tfrac12 JGJ,\qquad
 JGJ\equiv\int\dd^d x\dd^d y\,J(x)G(x,y)J(y).
\end{equation}
其中 \(\phi A\phi\) 和 \(J\phi\) 都包含积分。平移后关于 \(\eta\) 的高斯积分与零源相同，故其行列式因子在 \(\ZZ[J]/\ZZ[0]\) 中抵消，得到
\begin{equation}\label{eq:freeZ}
 Z[J]=e^{\frac12JGJ},\quad W[J]=\frac12JGJ,\quad
 \bar\phi(x)=\frac{\delta W}{\delta J(x)}=(GJ)(x),\quad G_{c,J}(x,y)=G(x,y).
\end{equation}
平均场因此满足 \(A\bar\phi=J\)。这和普通高斯变量的 \(a\bar q=J\) 是同一件事，只是数 \(a\) 换成了微分算符，逆数换成了积分核。

\physicsnote{\(G\) 有两种相容的读法：它衡量不同位置共同涨落的程度，也把一个源分布映射成平均场的变化。自由场中两种读法由同一个核实现。\(W\) 只有二次项，表示三阶及更高连通关联为零，绝不表示“两点关联也消失”。}
连续极限下，\(G(x,y)\) 可能在 \(x=y\) 奇异。上述式子应按分布理解，或先与平滑源积分；这里不把重合点发散当成一个有限的普通数。传播子的傅里叶表达与具体边界问题可在后续单元展开。

\newpage
\section{“响应”是哪一种响应？}
\question{既然 \(\delta\avg{\OO}/\delta J=G_c\)，是否已经算出了信号随时间传播？}
由式 \eqref{eq:covariance}，在小的欧氏源变化下，
\begin{equation}\label{eq:linear}
 \delta\avg{\OO(x)}_J
 =\int\dd^d y\,G_{c,J}(x,y)\,\delta J(y)+O(\delta J^2).
\end{equation}
这首先是在比较\textbf{不同欧氏权重下的期望值}。它告诉我们改变某处的权重，其他位置的平均值怎样变化；并不自动指定“先施加扰动、后观察结果”的因果实验。

\step{静态平衡：积分后的涨落控制易感率}
在有限温度下，欧氏时间为 \(0\leq\tau<\beta\)。令源不依赖欧氏时间，\(J(\tau,\mathbf y)=h\,j(\mathbf y)\)，则在零源处
\begin{equation}
 \left.\frac{\partial\avg{\OO(0,\mathbf x)}_h}{\partial h}\right|_{h=0}
 =\int_0^\beta\dd\tau\int\dd^{d-1}\mathbf y\,
 G_{c,0}((0,\mathbf x),(\tau,\mathbf y))\,j(\mathbf y).
\end{equation}
这才是某个指定外场形状对应的静态响应。对于量子系统，一般要保留欧氏时间积分；只有适当的对易条件或经典极限下，才能把它简化为 \(\beta\) 乘等时协方差；测量与外场共轭的同一个总观测量时，才成为方差。第 4 节的单变量模型正是最简单的静态例子。

\step{实时间：需要额外指定因果性}
若实验中扰动为
\(H'(t)=-\int\dd^{d-1}\mathbf y\,j(t,\mathbf y)\OO(t,\mathbf y)\)，对未扰动态取期望，可定义线性响应核
\begin{align}
 \delta\avg{\OO(t,\mathbf x)}
 &=\int\dd t'\dd^{d-1}\mathbf y\,
 \chi_R(t,\mathbf x;t',\mathbf y)\,j(t',\mathbf y),\notag\\
 \chi_R(t,\mathbf x;t',\mathbf y)
 &=i\,\theta(t-t')\avg{[\OO(t,\mathbf x),\OO(t',\mathbf y)]}_0.
\end{align}
这里的 \(+i\) 与上述 \(H'=-j\OO\) 及响应公式一起定义符号；也有文献把 \(-i\theta\avg{[\OO,\OO]}\) 定义为 \(G_R\)，此时本式的 \(\chi_R=-G_R\)。阶跃函数保证未来的源不会影响过去的读数。

\physicsnote{欧氏两点函数与实时间响应有联系，但要通过适当的解析延拓或实时间路径积分，并选择 retarded 条件。不能只把 \(\tau\) 改名为 \(t\)，也不能把欧氏 Green 核直接当作因果传播核。全息中选取内部正则解或入射解，正是这类额外物理信息的一部分。}

\newpage
\section{接到全息：边界条件为什么能生成关联函数？}
\question{为什么计算一个 bulk 在壳作用量，就能读出边界的均值和两点函数？}
先固定字典的约定：边界 QFT 采用式 \eqref{eq:definitions} 的 \(+\int J\OO\)；bulk 的渐近源数据识别为 \(\phi_{(0)}=J\)。欧氏字典示意为
\begin{equation}
 \ZZ_{\mathrm{QFT}}[J]
 =\ZZ_{\mathrm{bulk}}[\phi_{(0)}=J]
 \simeq e^{-I_{\mathrm{ren}}[\phi_{\mathrm{cl}}[J]]}.
\end{equation}
右侧近似要求适当的半经典极限，并假定当前由一个主要鞍点控制。\(\phi_{\mathrm{cl}}[J]\) 满足 bulk 方程、给定的渐近源以及内部条件；\(I_{\mathrm{ren}}\) 是加入边界项与反项后的重整化在壳作用量。

把右侧简记为 \(e^{-I_{\mathrm{ren}}[J]}\)，并采用本篇的零源归一化，便有
\begin{equation}\label{eq:holo}
 W[J]\simeq-I_{\mathrm{ren}}[J]+I_{\mathrm{ren}}[0],\qquad
 \avg{\OO(x)}_J\simeq-\frac{\delta I_{\mathrm{ren}}}{\delta J(x)},\qquad
 G_{c,J}(x,y)\simeq-\frac{\delta^2 I_{\mathrm{ren}}}{\delta J(x)\delta J(y)}.
\end{equation}
这里的负号由本页约定推出；比较文献时还要核对源的符号、Euclidean/Lorentzian 形式，以及径向法向和动量的定义，不能把负号当成独立于约定的字典。

\step{在壳以后为什么仍有非零变分？}
一般变分包含“运动方程乘 \(\delta\phi\)”与边界项。在壳只令前者为零；改变 \(J\) 时，边界数据也改变，边界项仍在。用与本页一致的定义写
\begin{equation}
 \delta I_{\mathrm{ren}}=\int\dd^d x\,\Pi_{\mathrm{ren}}(x)\,\delta J(x),
 \qquad \avg{\OO(x)}_J\simeq-\Pi_{\mathrm{ren}}(x).
\end{equation}
于是实际计算是：改变边界源，重新求满足内部条件的解，再观察其重整化边界动量怎样改变。通常渐近展开的一个系数指定源，另一个系数携带响应；两者之间的关系由 bulk 方程与内部条件共同决定。

\physicsnote{bulk 在半经典近似下使用经典解，并不表示边界 QFT 没有量子涨落：边界的连通二点函数已包含在 \(I_{\mathrm{ren}}[J]\) 对源的二阶依赖中。经典 bulk 与自由边界场也不是同义词；bulk 相互作用可以产生更高的边界连通关联。}
局部有限反项可以改变接触项，例如增加 \(\int J^2\) 会改变与 \(\delta(x-y)\) 成比例的二点项。因此裸的发散在壳作用量不能直接使用；讨论非重合点信息与重合点接触项时也须区分。

\newpage
\section{选读：能否用平均场代替源作为自变量？}
\question{\(W[J]\) 与有效作用量 \(\Gamma[\bar\phi]\) 分别回答什么问题？}
本节取源线性耦合到基本场 \(\phi\)。定义
\begin{equation}
 \bar\phi(x)=\frac{\delta W}{\delta J(x)},\qquad
 \Gamma[\bar\phi]=\intx J(x)\bar\phi(x)-W[J].
\end{equation}
假定在讨论的源附近可以把 \(\bar\phi[J]\) 局部反解为 \(J[\bar\phi]\)，右侧的 \(J\) 随之被消去。这个 Legendre 变换把“给定源，平均场是多少”改写成“要维持这个平均场，需要多大的源”。

现在直接变分：
\begin{align}
 \delta\Gamma
 &=\intx\bigl(\delta J\,\bar\phi+J\,\delta\bar\phi\bigr)
   -\intx\frac{\delta W}{\delta J}\delta J\notag\\
 &=\intx J\,\delta\bar\phi,
 \qquad \frac{\delta\Gamma}{\delta\bar\phi(x)}=J(x).
\end{align}
因此没有外源时，物理平均场满足 \(\delta\Gamma/\delta\bar\phi=0\)。这看起来像经典运动方程，但 \(\Gamma\) 一般已包含量子效应，不能直接用裸的 \(S_E\) 替代。

再用链式法则，对反函数关系求导：
\begin{equation}
 \int\dd^d z\,
 \frac{\delta^2\Gamma}{\delta\bar\phi(x)\delta\bar\phi(z)}
 \frac{\delta^2W}{\delta J(z)\delta J(y)}
 =\delta^{(d)}(x-y).
\end{equation}
在可逆的子空间上，这就是 \(\Gamma^{(2)}=G_c^{-1}\)。若有零模、规范冗余或不可反解之处，不能不加处理就写这个逆。

\physicsnote{\(W^{(2)}\) 描述平均场对源的灵敏度；\(\Gamma^{(2)}\) 描述把平均场推离原位置所需的“刚度”。第 4 节中二者分别为 \(1/a\) 与 \(a\)。更软的方向通常意味着更大的涨落与响应。}
自由场例子里 \(\bar\phi=GJ\)、\(J=A\bar\phi\)，所以
\begin{equation}
 \Gamma[\bar\phi]=\frac12\bar\phi A\bar\phi.
\end{equation}
这里恰好得到经典二次作用量；它是高斯模型的结果，不是相互作用理论的普遍恒等式。\(\Gamma\) 也称 1PI（one-particle-irreducible，一粒子不可约）有效作用量，其导数生成 1PI 顶角。全息中首先得到的是 \(W\simeq-I_{\mathrm{ren}}+\text{常数}\)，若要 \(\Gamma\)，还须按所选变量做相应的 Legendre 变换。

\newpage
\section{主动检验：用计算检查物理理解}
建议先遮住下面的核对线索，用 10--15 分钟完成前两题。第三题用于检验与全息的连接。

\textbf{1．均值不为零，如何分离背景与涨落？}
令 \(S(q)=a(q-b)^2/2\)、\(a>0\)，源仍为 \(+Jq\)，并仍用零源归一化。求 \(Z(J)\)、\(W(J)\)、\(\avg q_J\)、\(\avg{q^2}_J\)。判断：改变 \(b\) 后，零源的完整二点与连通二点是否相同？为什么？

\textbf{2．两个变量怎样互相响应？}
考虑
\begin{equation}
 S(q_1,q_2)=\frac a2(q_1^2+q_2^2)+c q_1q_2,
 \qquad a>|c|,
\end{equation}
源为 \(J_1q_1+J_2q_2\)。不必逐个积分，利用二次型求
\(\partial\avg{q_1}/\partial J_2\)。当 \(c>0\) 时为什么是负的？当 \(c\to a^-\) 时，哪一个组合变软，响应为何增大？

\textbf{3．只知道边界源，为什么还不够？}
用自己的话解释：求全息响应时，为什么除了边界源还要给内部条件？若把欧氏计算得到的二点核直接叫作实时间因果响应，缺少了什么？

\medskip\hrule\medskip
\textbf{核对线索}
\begin{enumerate}
\item 平移 \(q=b+\eta\)：\(Z=e^{bJ+J^2/(2a)}\)、\(W=bJ+J^2/(2a)\)。均值为 \(b+J/a\)，完整二点为 \((b+J/a)^2+1/a\)，连通二点为 \(1/a\)。背景均值与涨落宽度是不同信息。
\item 矩阵 \(A=\left(\begin{smallmatrix}a&c\\c&a\end{smallmatrix}\right)\)，故交叉响应为
\(-c/(a^2-c^2)\)。\(c>0\) 惩罚同号构型；\(c\to a^-\) 时反向组合 \(q_1-q_2\) 的刚度 \(a-c\) 趋零。端点处不再正定，本高斯积分失效。
\item 边界源只指定问题的一部分；内部条件选择相应的解与物理状态。实时间还要指定因果处方，例如 retarded 条件，不能只改时间符号。
\end{enumerate}

\textbf{本单元完成标准。}能够不看材料推导 \(W\) 的前两阶导数，解释归一化为何产生减项；能算完高斯例子并预测均值、涨落与响应怎样变化；能说明全息中“源作为边界数据、在壳边界变分读取响应”的逻辑。第 8 节不作为本次必须掌握的门槛。

\textbf{延伸阅读与约定对照}
\begin{itemize}
\item D. Tong, \textit{Statistical Field Theory}, \href{https://davidtong.org/teaching/statistical-field-theory/sfthtml/S2}{第 2 章，\S2.2--2.2.2}：源、连通关联与 Green 函数。其统计物理记号含 \(\beta\)，与本文比较时先换算源。
\item D. Skinner, \textit{Advanced Quantum Field Theory}, \href{https://www.damtp.cam.ac.uk/user/dbs26/AQFT/chap2.pdf}{第 2 章}：零维场论的高斯例子及进一步展开。
\item J. McGreevy, \textit{Holographic duality with a view toward many-body physics}, \href{https://arxiv.org/abs/0909.0518}{arXiv:0909.0518，第 4 节}：全息生成泛函与源、响应；特别参看 \S4.4。比较公式前请核对各自符号约定。
\end{itemize}
\end{document}
